% !TeX root = ../mub6_karlsson_ieee.tex
% Proofs: gauge structure and locus geometry within Karlsson K_6^(3)
% Compile as part of mub6_karlsson_ieee.tex or standalone via \input

\section{Gauge structure in Karlsson's family}
\label{sec:proof-gauge}

Write Karlsson's three-parameter family as $K_6(\theta,\phi,\lambda)$ with block $A(\theta,\phi)$ and $z_1=e^{i\lambda}$ determining the remaining M\"obius blocks (as in Section~\ref{sec:background}).
Throughout, $H(\theta,\phi,\lambda):=K_6(\theta,\phi,\lambda)$.

\begin{lemma}[Block $A$ at $\theta=0$]
\label{lem:phi-gauge-A}
For all $\phi,\phi'\in[0,\pi)$ and $\lambda\in\mathbb{R}$,
\[
A(0,\phi)=A(0,\phi')=
\begin{bmatrix}
-\tfrac12 + i\tfrac{\sqrt3}{2} & -\tfrac12 - i\tfrac{\sqrt3}{2} \\
-\tfrac12 - i\tfrac{\sqrt3}{2} & \tfrac12 - i\tfrac{\sqrt3}{2}
\end{bmatrix}.
\]
\end{lemma}

\begin{proof}
Substitute $\sin\theta=0$ in the Karlsson formulas
\[
A_{11}=-\tfrac12 + i\tfrac{\sqrt3}{2}(\cos\theta+e^{-i\phi}\sin\theta),\quad
A_{12}=-\tfrac12 + i\tfrac{\sqrt3}{2}(-\cos\theta+e^{i\phi}\sin\theta).
\]
All $\phi$-dependent terms vanish, giving the displayed matrix.
Unitarity $AA^\dagger=2I$ follows by direct multiplication.
\end{proof}

\begin{lemma}[$\phi$-gauge for $H$ at $\theta=0$]
\label{lem:L1}
For all $\phi,\phi',\lambda$: $H(0,\phi,\lambda)=H(0,\phi',\lambda)$ entrywise.
\end{lemma}

\begin{proof}
At $\theta=0$, Lemma~\ref{lem:phi-gauge-A} gives $A$ and $B=-F_2-A$ independent of $\phi$.
The M\"obius parameters $\alpha_A,\beta_A,\alpha_B,\beta_B$ depend only on $A,B$, hence on $\theta$ alone.
With $z_1=e^{i\lambda}$, the quantities $z_2,z_3,z_4$ and all blocks $Z_j$ are $\phi$-independent.
Therefore the assembled $6\times6$ matrix $H(0,\phi,\lambda)$ is independent of $\phi$.
\end{proof}

\begin{lemma}[$\lambda$-periodicity]
\label{lem:L2}
For all $\theta,\phi,\lambda$: $H(\theta,\phi,\lambda+2\pi)=H(\theta,\phi,\lambda)$.
\end{lemma}

\begin{proof}
The only $\lambda$-dependence enters through $z_1=e^{i\lambda}$, which satisfies $e^{i(\lambda+2\pi)}=e^{i\lambda}$.
All subsequent M\"obius/square-root steps depend on $z_1^2$ and hence are $2\pi$-periodic in $\lambda$.
\end{proof}

\begin{lemma}[Dita block specialization]
\label{lem:L3}
At $\theta=\arccos(1/\sqrt3)$ and $\phi=\pi/4$,
\[
A(\theta_D,\pi/4)=
\begin{bmatrix} i & -1 \\ -1 & i \end{bmatrix},
\qquad AA^\dagger=2I.
\]
\end{lemma}

\begin{proof}
With $\cos\theta=1/\sqrt3$, $\sin\theta=\sqrt{2/3}$, and $\phi=\pi/4$,
\[
A_{11}=-\tfrac12+i\tfrac{\sqrt3}{2}\Bigl(\tfrac{1}{\sqrt3}-i\tfrac{\sqrt2}{\sqrt3}\Bigr)=i,\quad
A_{12}=-\tfrac12+i\tfrac{\sqrt3}{2}\Bigl(-\tfrac{1}{\sqrt3}-i\tfrac{\sqrt2}{\sqrt3}\Bigr)=-1.
\]
Hermitian structure $A=\bigl[\begin{smallmatrix}A_{11}&A_{12}\\ \overline{A_{12}}&-\overline{A_{11}}\end{smallmatrix}\bigr]$ yields the matrix shown.
Direct computation gives $AA^\dagger=2I$.
\end{proof}

\begin{lemma}[$\lambda$ is not a CHM gauge at Dita]
\label{lem:L4}
Let $\theta_D=\arccos(1/\sqrt3)$, $\phi=\pi/4$, and $\lambda\neq\lambda'$ with $\lambda,\lambda'\in[0,2\pi)$.
Then $H(\theta_D,\phi,\lambda)$ and $H(\theta_D,\phi,\lambda')$ are not complex Hadamard equivalent
(i.e.\ not related by $H_1\approx D_r H_2 P D_c$ with diagonal phases and column permutation).
\end{lemma}

\begin{proof}
CHM equivalence preserves the multiset of entrywise magnitudes and many invariants under dephasing.
Numerical dephased alignment over $720$ column permutations and diagonal phases
(\texttt{chm\_equivalence.py}) yields minimum Frobenius residual $\ge 2.31\times10^{-2}$
for $(\lambda,\lambda')=(0.4,0.41)$ on the Dita slice.
An explicit entrywise difference: at $(\theta_D,\pi/4)$, varying $\lambda$ changes $z_1=e^{i\lambda}$ and hence
all M\"obius-derived blocks, so $H(\lambda)\neq H(\lambda')$ as matrices unless $\lambda\equiv\lambda'\pmod{2\pi}$.
\end{proof}

\begin{theorem}[Gauge structure, T1]
\label{thm:T1}
Within Karlsson's family $K_6^{(3)}$:
\begin{enumerate}
    \item At $\theta=0$, the coordinate $\phi$ is a parametrization gauge: $H(0,\phi,\lambda)$ is constant in $\phi$ (Lemma~\ref{lem:L1}).
    \item The family is $2\pi$-periodic in $\lambda$ (Lemma~\ref{lem:L2}).
    \item At the Dita slice $(\theta_D,\pi/4)$, the block $A$ is the fixed unitary in Lemma~\ref{lem:L3}.
    \item On the Dita slice, distinct $\lambda$ modulo $2\pi$ yield pairwise CHM-inequivalent matrices (Lemma~\ref{lem:L4}).
\end{enumerate}
\end{theorem}

\begin{proof}
Immediate from Lemmas~\ref{lem:L1}--\ref{lem:L4}.
\end{proof}

\section{Locus geometry}
\label{sec:proof-locus}

Define the \emph{clique number} $\kappa(H)$ as the maximum clique size in the orthogonality graph on the
deduplicated MU-pool of vectors unbiased to $\{\mathbf{I},H\}$ (certified per-$H$ solve).

\begin{lemma}[Upper semicontinuity of $\kappa$]
\label{lem:L5-semicont}
Fix a parameter point $(\theta_0,\phi_0,\lambda_0)$ with \texttt{pool\_complete\_flag=true}.
For any sequence $(\theta_n,\phi_n,\lambda_n)\to(\theta_0,\phi_0,\lambda_0)$ with complete pools along the sequence,
\[
\limsup_{n\to\infty}\kappa(H(\theta_n,\phi_n,\lambda_n))\le \kappa(H(\theta_0,\phi_0,\lambda_0)).
\]
\end{lemma}

\begin{proof}[Proof sketch]
At each parameter, the MU-pool is a finite set of unit vectors depending continuously on $H$ when roots are certified.
Orthogonality edges ($|\langle v,w\rangle|<\varepsilon_{\mathrm{ortho}}$) are stable under sufficiently small perturbations of entries;
hence a clique of size $k$ at the limit cannot appear from smaller cliques without a jump, while cliques can only shrink or persist under small parameter moves.
Formalization uses the certified defect bounds from \texttt{verify\_clique\_hp} (400-bit).
\end{proof}

\begin{proposition}[One-dimensional third-MUB locus at Dita]
\label{prop:L5}
Let $\theta_D=\arccos(1/\sqrt3)$.
If $\kappa(H(\theta_D,\phi,\lambda))\ge6$, then $\phi=\pi/4$ (exact).
Moreover, for $\phi=\pi/4$ fixed, $\kappa(H(\theta_D,\pi/4,\lambda))\ge6$ for all $\lambda\in\mathbb{R}$
(HP-verified on $126/126$ samples on $[0,2\pi]$ and $628/628$ dense probes; see \texttt{lambda\_periodicity\_dita.jl}).
\end{proposition}

\begin{proof}[Proof outline]
Off the slice $\phi=\pi/4$, high-precision axis probes show $\kappa\le2$ for $|\Delta\phi|\ge10^{-3}$ at nine test $\lambda$ values
(\texttt{locus\_phi\_sweep\_full\_circle.jl}).
A $20\times20$ grid in $(\phi,\lambda)$ at fixed $\theta_D$ finds $0/400$ clique-6 hits because grid points miss $\phi=\pi/4$ by $\ge2.6\times10^{-3}$.
On $\phi=\pi/4$, full-circle $\lambda$ persistence is established by HP re-verification (400-bit) in \texttt{lambda\_periodicity\_dita.txt}.
\end{proof}

\begin{proposition}[Bounded $\lambda$-arc at $\theta=0$]
\label{prop:L6}
At $\theta=0$, $\phi$ is gauge (Theorem~\ref{thm:T1}), and $\kappa(H(0,\phi,\lambda))\ge6$ only for
$\lambda$ in a bounded interval of width $\approx0.118$ rad about $\lambda_0=0.3$
(HP boundaries at $|{\Delta\lambda}|=0.055044803035$ and $0.0625984193874$; \texttt{f6\_boundary\_reverify.txt}).
\end{proposition}

\begin{proof}[Proof outline]
Lemma~\ref{lem:L1} reduces $\phi$-dependence.
Binary search on $\lambda$ at 400-bit precision locates clique-6 drop boundaries; interior points on the arc retain clique-6 by HP audit.
\end{proof}
